\begin{figure}[htbp]
\centering
\begin{tikzpicture}[x=1cm,y=1cm]
\begin{scope}
\draw[thin] (0,2.4)--(0,0)--(6,0);
% Original printed PDF vector path.
\path[draw=plotgray,line width=.5pt,fill=gray!55] (0.50002,0.20001) .. controls (0.50002,0.20001) and (0.57375,0.20002) ..(0.60206,0.20003) .. controls (0.63038,0.20004) and (0.67579,0.20006) ..(0.70410,0.20011) .. controls (0.73242,0.20016) and (0.77783,0.20023) ..(0.80614,0.20038) .. controls (0.83446,0.20053) and (0.87987,0.20077) ..(0.90818,0.20119) .. controls (0.93650,0.20161) and (0.98191,0.20233) ..(1.01022,0.20343) .. controls (1.03854,0.20453) and (1.08395,0.20650) ..(1.11226,0.20914) .. controls (1.14058,0.21177) and (1.18599,0.21665) ..(1.21430,0.22242) .. controls (1.24262,0.22819) and (1.28803,0.23916) ..(1.31634,0.25072) .. controls (1.34466,0.26228) and (1.39007,0.28460) ..(1.41838,0.30574) .. controls (1.44670,0.32689) and (1.49211,0.36789) ..(1.52042,0.40314) .. controls (1.54874,0.43840) and (1.59415,0.50656) ..(1.62246,0.55985) .. controls (1.65078,0.61314) and (1.69619,0.71460) ..(1.72450,0.78722) .. controls (1.75282,0.85985) and (1.79823,0.99485) ..(1.82654,1.08328) .. controls (1.85486,1.17171) and (1.90027,1.33005) ..(1.92858,1.42456) .. controls (1.95690,1.51907) and (2.00231,1.67880) ..(2.03062,1.76446) .. controls (2.05894,1.85013) and (2.10435,1.98167) ..(2.13266,2.04198) .. controls (2.16098,2.10230) and (2.20639,2.17726) ..(2.23470,2.19918) .. controls (2.26302,2.22110) and (2.30842,2.22128) ..(2.33674,2.19999) .. controls (2.36506,2.17870) and (2.41046,2.10481) ..(2.43878,2.04572) .. controls (2.46710,1.98664) and (2.51250,1.85714) ..(2.54082,1.77415) .. controls (2.56914,1.69115) and (2.61454,1.53630) ..(2.64286,1.44756) .. controls (2.67118,1.35882) and (2.71658,1.21147) ..(2.74490,1.13457) .. controls (2.77322,1.05767) and (2.81862,0.94488) ..(2.84694,0.89332) .. controls (2.87526,0.84175) and (2.92066,0.78104) ..(2.94898,0.76295) .. controls (2.97730,0.74487) and (3.02270,0.74495) ..(3.05102,0.76298) .. controls (3.07933,0.78100) and (3.12474,0.84143) ..(3.15306,0.89289) .. controls (3.18137,0.94435) and (3.22678,1.05703) ..(3.25510,1.13386) .. controls (3.28341,1.21069) and (3.32882,1.35787) ..(3.35714,1.44660) .. controls (3.38545,1.53533) and (3.43086,1.69029) ..(3.45918,1.77336) .. controls (3.48749,1.85643) and (3.53290,1.98615) ..(3.56122,2.04531) .. controls (3.58953,2.10447) and (3.63494,2.17840) ..(3.66326,2.19977) .. controls (3.69157,2.22114) and (3.73698,2.22117) ..(3.76530,2.19936) .. controls (3.79361,2.17756) and (3.83902,2.10282) ..(3.86734,2.04259) .. controls (3.89565,1.98236) and (3.94106,1.85089) ..(3.96938,1.76525) .. controls (3.99769,1.67961) and (4.04310,1.51986) ..(4.07142,1.42535) .. controls (4.09973,1.33083) and (4.14514,1.17253) ..(4.17346,1.08408) .. controls (4.20177,0.99563) and (4.24718,0.86053) ..(4.27550,0.78785) .. controls (4.30381,0.71517) and (4.34922,0.61360) ..(4.37754,0.56026) .. controls (4.40585,0.50691) and (4.45126,0.43863) ..(4.47958,0.40335) .. controls (4.50789,0.36806) and (4.55330,0.32709) ..(4.58162,0.30592) .. controls (4.60993,0.28476) and (4.65534,0.26240) ..(4.68366,0.25082) .. controls (4.71197,0.23924) and (4.75738,0.22825) ..(4.78570,0.22247) .. controls (4.81401,0.21669) and (4.85942,0.21180) ..(4.88774,0.20916) .. controls (4.91605,0.20652) and (4.96146,0.20454) ..(4.98978,0.20344) .. controls (5.01809,0.20233) and (5.06350,0.20162) ..(5.09182,0.20119) .. controls (5.12013,0.20077) and (5.16554,0.20053) ..(5.19386,0.20038) .. controls (5.22217,0.20023) and (5.26758,0.20016) ..(5.29589,0.20011) .. controls (5.32421,0.20006) and (5.36962,0.20005) ..(5.39793,0.20003) .. controls (5.42625,0.20002) and (5.49998,0.20001) ..(5.49998,0.20001);

\node[below] at (3,-.05){$x$};
\node[rotate=90] at (-.42,1.2){$p(x)$};
\end{scope}
\begin{scope}[xshift=7.4cm]
\draw[thin] (0,2.4)--(0,0)--(6,0);
% Original printed PDF vector path.
\path[draw=plotgray,line width=.5pt,fill=gray!55] (0.50002,0.20002) .. controls (0.50002,0.20002) and (0.57375,0.20003) ..(0.60206,0.20004) .. controls (0.63038,0.20006) and (0.67579,0.20007) ..(0.70410,0.20012) .. controls (0.73242,0.20017) and (0.77783,0.20024) ..(0.80614,0.20039) .. controls (0.83446,0.20054) and (0.87987,0.20077) ..(0.90818,0.20120) .. controls (0.93650,0.20162) and (0.98191,0.20234) ..(1.01022,0.20344) .. controls (1.03854,0.20455) and (1.08395,0.20651) ..(1.11226,0.20915) .. controls (1.14058,0.21179) and (1.18599,0.21667) ..(1.21430,0.22244) .. controls (1.24262,0.22821) and (1.28803,0.23919) ..(1.31634,0.25075) .. controls (1.34466,0.26232) and (1.39007,0.28464) ..(1.41838,0.30580) .. controls (1.44670,0.32696) and (1.49211,0.36797) ..(1.52042,0.40325) .. controls (1.54874,0.43852) and (1.59415,0.50671) ..(1.62246,0.56003) .. controls (1.65078,0.61334) and (1.69619,0.71485) ..(1.72450,0.78750) .. controls (1.75282,0.86016) and (1.79823,0.99522) ..(1.82654,1.08369) .. controls (1.85486,1.17216) and (1.90026,1.33057) ..(1.92858,1.42513) .. controls (1.95690,1.51968) and (2.00231,1.67948) ..(2.03062,1.76517) .. controls (2.05894,1.85087) and (2.10434,1.98245) ..(2.13266,2.04276) .. controls (2.16098,2.10307) and (2.20639,2.17801) ..(2.23470,2.19983) .. controls (2.26302,2.22165) and (2.30842,2.22165) ..(2.33674,2.20003) .. controls (2.36505,2.17841) and (2.41046,2.10395) ..(2.43878,2.04402) .. controls (2.46710,1.98408) and (2.51250,1.85304) ..(2.54082,1.76804) .. controls (2.56914,1.68303) and (2.61454,1.52445) ..(2.64286,1.43134) .. controls (2.67117,1.33823) and (2.71658,1.18258) ..(2.74490,1.09698) .. controls (2.77322,1.01137) and (2.81862,0.88177) ..(2.84694,0.81434) .. controls (2.87525,0.74691) and (2.92066,0.65549) ..(2.94898,0.61096) .. controls (2.97730,0.56643) and (3.02270,0.51535) ..(3.05102,0.49340) .. controls (3.07933,0.47145) and (3.12474,0.45577) ..(3.15306,0.45276) .. controls (3.18137,0.44975) and (3.22678,0.46116) ..(3.25510,0.47170) .. controls (3.28341,0.48223) and (3.32882,0.51080) ..(3.35714,0.52866) .. controls (3.38545,0.54652) and (3.43086,0.58161) ..(3.45918,0.60040) .. controls (3.48749,0.61920) and (3.53290,0.65017) ..(3.56122,0.66414) .. controls (3.58953,0.67811) and (3.63494,0.69606) ..(3.66326,0.70108) .. controls (3.69157,0.70611) and (3.73698,0.70594) ..(3.76530,0.70038) .. controls (3.79361,0.69481) and (3.83902,0.67606) ..(3.86734,0.66096) .. controls (3.89565,0.64586) and (3.94106,0.61296) ..(3.96938,0.59153) .. controls (3.99769,0.57010) and (4.04310,0.53014) ..(4.07142,0.50650) .. controls (4.09973,0.48286) and (4.14514,0.44326) ..(4.17346,0.42114) .. controls (4.20177,0.39901) and (4.24718,0.36522) ..(4.27550,0.34704) .. controls (4.30381,0.32886) and (4.34922,0.30346) ..(4.37754,0.29012) .. controls (4.40585,0.27677) and (4.45126,0.25970) ..(4.47958,0.25087) .. controls (4.50789,0.24205) and (4.55330,0.23180) ..(4.58162,0.22650) .. controls (4.60993,0.22121) and (4.65534,0.21562) ..(4.68366,0.21272) .. controls (4.71197,0.20983) and (4.75738,0.20708) ..(4.78569,0.20563) .. controls (4.81401,0.20419) and (4.85942,0.20296) ..(4.88774,0.20230) .. controls (4.91605,0.20164) and (4.96146,0.20115) ..(4.98978,0.20087) .. controls (5.01809,0.20059) and (5.06350,0.20041) ..(5.09181,0.20031) .. controls (5.12013,0.20020) and (5.16554,0.20014) ..(5.19386,0.20011) .. controls (5.22217,0.20007) and (5.26758,0.20005) ..(5.29589,0.20004) .. controls (5.32421,0.20003) and (5.36962,0.20002) ..(5.39793,0.20002) .. controls (5.42625,0.20001) and (5.49997,0.20001) ..(5.49997,0.20001);

\node[below] at (3,-.05){$x$};
\node[rotate=90] at (-.42,1.2){$p'(x)$};
\end{scope}
\end{tikzpicture}
\caption{$P_X$ 变为 $P'_X$ 的一种情形：这种变化提示，$P_Y$ 已改变，而 $P_{X\mid Y}$ 保持不变。当已知二元变量 $Y$ 是 $X$ 的原因时，观测到 $P_X$ 为两个高斯分布的混合，会使我们有理由认为，两个峰分别对应不同的标签 $y=0,1$。于是，$Y$ 对 $X$ 的影响只是平移高斯分布的均值，这就构成加性噪声模型，见第 4.1.4 节。这无疑为联合分布提供了简单解释。如果进一步观测到，混合权重在两个数据集之间发生了变化，那么这一变化就很可能是由 $P_Y$ 的变化造成的。}
\label{fig:ch05-changing-mixture}
\end{figure}
